\documentclass[10pt,a4paper]{amsart}
\usepackage[margin=19mm]{geometry}
\usepackage{amsmath,amssymb,mathtools}
\usepackage[scr=rsfs,cal=euler]{mathalfa}
\usepackage{xfrac}
\usepackage[unicode,hidelinks,bookmarks=false]{hyperref}
\newcommand{\ie}{i.e.\ }
\newcommand{\cf}{cf.\ }
\newcommand{\resp}{respectively\ }
\newcommand{\Subsection}{\subsection}
\providecommand{\nobreakdash}{\nobreak\hbox{-}}
\setlength{\emergencystretch}{2em}
\multlinegap0pt
\usepackage[shortlabels]{enumitem}
\usepackage{amssymb}
\usepackage{float}
\usepackage{xcolor}
\usepackage{mathtools}
\usepackage{csquotes}
\usepackage[normalem]{ulem}
\newtheorem{theorem}{Theorem}
\usepackage{cleveref}
\crefname{theorem}{Theorem}{Theorems}
\title{Exterior differential systems on Lie algebroids and the invariant inverse problem of the calculus of variations}
\author{Tom Mestdag, Kenzo Yasaka}
\date{Source: arXiv:2507.14678v1 (CC BY 4.0)}
\begin{document}
\maketitle

\noindent\emph{Source and licence.} Exterior differential systems on Lie algebroids and the invariant inverse problem of the calculus of variations. Version arXiv:2507.14678v1 (CC BY 4.0), distributed by arXiv under the Creative Commons Attribution 4.0 International licence, http://creativecommons.org/licenses/by/4.0/, as recorded in the arXiv licence metadata for this exact version and shown on the abstract page https://arxiv.org/abs/2507.14678v1. Full licence notice: \texttt{licenses/arxiv-2507.14678-CC-BY-4.0.txt}.\par\smallskip

\noindent\textit{Editorial scope.} Counted unit: the article's theorem thm:helmholtzintmfd (source0.tex lines 2183--2189). The article's complete original proof of it is delivered with it (source0.tex lines 2190--2237, one \texttt{proof} environment of 48 lines). No mathematics is written, repaired or re-derived: the statement and the proof are the article's own lines, in the article's own notation, and no retained line is substituted or re-flowed.  Retained uncounted as the source-local context the proof needs: lines 2024--2026.  Nothing was omitted from any retained span, so the package carries no omission record.  No retained line contains a citation, so the wrapper carries no bibliography and no bibliography entry.  No retained line contains a figure environment, an image-inclusion command or drawing code, so no image is delivered and the package declares no asset dependency.  Editorial formatting only, in addition: an AMS \texttt{amsart} preamble that restates the article's own declarations and macros used by the retained lines, namely the environments \texttt{theorem} on separate counters, together with the standard packages those lines need.  The retained environments print in this document as Theorem 1 in order of appearance, whereas the article prints the same items with the numbering of its own full document.  The complete native source of the article as distributed in the arXiv e-print for this exact version is bundled unchanged as \texttt{sources/181871/source0.tex}.
\begin{equation}\label{eqn:redundanthelmholtz1}
    \psi_l^k \frac{\partial k_{ij}}{\partial w^k}+\frac{1}{2}\left(k_{ki}C_{jl}^k+k_{kj}C_{il}^k-k_{kl}\frac{\partial^2 \gamma^k}{\partial w^i \partial w^j}\right)=\psi_j^k \frac{\partial k_{il}}{\partial w^k}+\frac{1}{2}\left(k_{ik} C_{lj}^k+k_{kl} C_{ij}^k-k_{kj}\frac{\partial^2 \gamma^k}{\partial w^i \partial w^l}\right),
\end{equation}

\begin{theorem}\label{thm:helmholtzintmfd}
  Suppose that $\langle \Sigma^0\rangle_{\mathrm{alg}}$ is a differential ideal. Let $\mathcal{I}=\langle \sigma_{ij}\rangle$. A solution $(k_{ij})$ to the reduced Helmholtz conditions exists if and only if there exists an integral manifold $i\colon M\to \overline{M}$,
\begin{equation*}
    (t,w)\longmapsto \left(t,w,s_{ij}=\bar{s}_{ij}(t,w),P_{ijk}=\overline{P}_{ijk}(t,w),Q_{ijk}=\overline{Q}_{ijk}(t,w)\right)
\end{equation*}
of $\mathcal{I}$ on  the IP Lie algebroid, with $\det(\bar{s}_{ij})\neq 0$.
\end{theorem}

\begin{proof}
We first do some preparatory computations. Note that
\begin{equation*}
    I^{\ast}(s_{ij})=\gamma(\bar{s}_{ij})+\frac{\partial \bar{s}_{ij}}{\partial w^l}\Psi^l+\psi_l^k \frac{\partial \bar{s}_{ij}}{\partial w^k}\Theta^l.
\end{equation*}
Hence
\begin{multline*}
    I^{\ast}(\sigma_{ij})=\left[\gamma(\bar{s}_{ij})-\bar{s}_{ki}\lambda_j^k-\bar{s}_{kj}\lambda_i^k\right]\Gamma^0+\left[\overline{P}_{ijl}+\frac{\partial \bar{s}_{ij}}{\partial w^l}\right]\Psi^l\\+\left[\psi_l^k \frac{\partial \bar{s}_{ij}}{\partial w^k}+\frac{1}{2}\left(\bar{s}_{ki}C_{jl}^k+\bar{s}_{kj}C_{il}^k-\bar{s}_{kl}\frac{\partial^2 \gamma^k}{\partial w^i \partial w^j}\right)+\overline{Q}_{ijl}\right]\Theta^l.
\end{multline*}
Note that $I^{\ast}(\sigma_{ij})=0$  if and only if the following equations hold:
\begin{enumerate}[label=(\alph*)]
   % \item $\det(\bar{s}_{ij})\neq 0$.
    \item $\gamma(\bar{s}_{ij})-\bar{s}_{ki}\lambda_j^k-\bar{s}_{kj}\lambda_i^k=0$,
    \item $\overline{P}_{ijl}+\frac{\partial \bar{s}_{ij}}{\partial w^l}=0$ for all $l\in \{1,\dots,n\}$,
    \item $\psi_l^k \frac{\partial \bar{s}_{ij}}{\partial w^k}+\frac{1}{2}\left(\bar{s}_{ki}C_{jl}^k+\bar{s}_{kj}C_{il}^k-\bar{s}_{kl}\frac{\partial^2 \gamma^k}{\partial w^i \partial w^j}\right)+\overline{Q}_{ijl}=0$ for all $l\in \{1,\dots,n\}$.
\end{enumerate}
Suppose first that a solution to the reduced Helmholtz conditions exists. Denote this solution by $(k_{ij})=(\bar{s}_{ij})$. Then it follows directly that (a) holds. We now claim that
\begin{equation*}
    \overline{P}_{ijl}=-\frac{\partial \bar{s}_{ij}}{\partial w^l},\quad \overline{Q}_{ijl}=-\psi_l^k \frac{\partial \bar{s}_{ij}}{\partial w^k}-\frac{1}{2}\left(\bar{s}_{ki}C_{jl}^k+\bar{s}_{kj}C_{il}^k-\bar{s}_{kl}\frac{\partial^2 \gamma^k}{\partial w^i \partial w^j}\right),
\end{equation*}
is a solution to equations (b) and (c). The only requirement that these functions must satisfy is that $\overline{P}_{ijl}$ and $\overline{Q}_{ijl}$ be totally symmetric. By construction, they are both symmetric in $ij$. Furthermore, note that
\begin{equation*}
    \overline{P}_{ijl}=-\frac{\partial \bar{s}_{ij}}{\partial w^l}=-\frac{\partial \bar{s}_{lj}}{\partial w^i}=\overline{P}_{lji}.
\end{equation*}
As such, $\overline{P}_{ijl}$ is also symmetric in $il$. Hence $\overline{P}_{ijl}$ is totally symmetric. Similarly, $\overline{Q}_{ijl}$ is symmetric in the indices $jl$, because by one of the redundant Helmholtz conditions (specifically, \Cref{eqn:redundanthelmholtz1}), and the symmetry of $\bar{s}_{ij}$, one gets
\begin{equation*}
    \overline{Q}_{ijl}=-\psi_j^k \frac{\partial \bar{s}_{il}}{\partial w^k}-\frac{1}{2}\left(\bar{s}_{ki}C_{lj}^k+\bar{s}_{kl}C_{ij}^k-\bar{s}_{kj}\frac{\partial^2 \gamma^k}{\partial w^i \partial w^l}\right)=\overline{Q}_{ilj},
\end{equation*}
which implies that $\overline{Q}_{ijl}$ is totally symmetric.

Conversely, suppose that we are given a solution $(\bar{s}_{ij},\overline{P}_{ijk}, \overline{Q}_{ijk})$ to the system of equations consisting of (a), (b), (c) and $\det(\bar{s}_{ij})\neq 0$. We claim that $(\bar{s}_{ij})$ satisfies the reduced Helmholtz conditions. Note that if $\langle \Sigma^0\rangle_{\mathrm{alg}}$ is a differential ideal, the condition
\begin{equation*}
    \bar{s}_{ki}\phi_j^k=\bar{s}_{kj}\phi_i^k
\end{equation*}
holds by assumption. It therefore remains to prove that
\begin{equation}\label{eqn:helmholtzDV}
    \frac{\partial \bar{s}_{jk}}{\partial w^i}=\frac{\partial \bar{s}_{ik}}{\partial w^j}.
\end{equation}
But, by equation (b), we obtain
\begin{equation*}
    \overline{P}_{jki}+\frac{\partial \bar{s}_{jk}}{\partial w^i}=0,
\end{equation*}
and
\begin{equation*}
    \overline{P}_{ikj}+\frac{\partial \bar{s}_{ik}}{\partial w^j}=0.
\end{equation*}
Since $\overline{P}_{ijk}$ is totally symmetric, it follows that $\overline{P}_{ikj}=\overline{P}_{jki}$. Hence, $\bar{s}_{ij}$ must satisfy \Cref{eqn:helmholtzDV}, as required.
\end{proof}

\end{document}
